% !TEX root = ../main.tex
% §2 Related Work. Every number in Table~\ref{tab:related} was checked against the cited paper:
% SPADE / Rippel / PaDiM on MVTec AD from Table 2 of Defard et al.; PatchCore on MVTec AD from its abstract;
% VisA rows from Tables 2 and 6 of Zou et al.; "Ours" from results/stage6_main/summary.csv (12 classes).
\section{Related Work}
\label{sec:related}

\paragraph{Benchmarks and methods.}
Our task is in the image data category, where industrial inspection is usually studied on MVTec AD
\citep{bergmann2019mvtec}, where image-level AUROC is close to saturated (Table~\ref{tab:related}).
VisA \citep{zou2022spot} adds complex structures and freely placed instances, and the same PatchCore drops
there from 99.8 to 92.4 AUROC, so it still separates methods.
Among the families reviewed by \citet{ruff2021unifying}, the strongest normal-only detectors on such data
pair ImageNet-pretrained features with a simple model: a global Gaussian \citep{rippel2021modeling}, a
Gaussian per location (PaDiM, \citealp{defard2021padim}), or $k$NN to normal patches (SPADE,
\citealp{cohen2020spade}; PatchCore, \citealp{roth2022towards}).

\begin{table*}[t]
  \centering
  \small
  \setlength{\tabcolsep}{5pt}
  \begin{tabular}{@{}lllccc@{}}
    \toprule
    Work & Data & Method (family), backbone & Paradigm & AUROC & AU-PR \\
    \midrule
    \citet{cohen2020spade}     & MVTec AD & SPADE (patch $k$NN), WRN50 & normal only & 85.5$^\dagger$ & -- \\
    \citet{rippel2021modeling} & MVTec AD & global Gaussian, EfficientNet-B4 & normal only & 95.8 & -- \\
    \citet{defard2021padim}    & MVTec AD & PaDiM (patch Gaussian), WRN50 & normal only & 95.3 & -- \\
    \citet{roth2022towards}    & MVTec AD & PatchCore (patch $k$NN) & normal only & up to 99.6 & -- \\
    \citet{zou2022spot}        & VisA & PaDiM (patch Gaussian), ResNet-50 & normal only & 87.8 & 88.2 \\
    \citet{zou2022spot}        & VisA & PatchCore (patch $k$NN), WRN50 & normal only & 92.4 & 93.3 \\
    \midrule
    Ours, Tier 2 & VisA$^\ast$ & PaDiM (patch Gaussian), WRN50 & normal + val.\ $\tau$ & 82.8\,$\pm$\,0.8 & 80.1\,$\pm$\,0.7 \\
    Ours, Tier 3 & VisA$^\ast$ & PatchCore (patch $k$NN), WRN50 & normal + val.\ $\tau$ & 90.2\,$\pm$\,0.2 & 88.5\,$\pm$\,0.2 \\
    \bottomrule
  \end{tabular}
  \caption{Image-level results (\%) of prior work and ours. ``Normal + val.\ $\tau$'': labelled val
    anomalies only choose thresholds and hyperparameters. $^\dagger$As reported by
    \citet{defard2021padim}. $^\ast$Our split (\S\ref{sec:data}): context, not a direct comparison;
    mean\,$\pm$\,std over 3 seeds, 12 objects.}
  \label{tab:related}
\end{table*}

\paragraph{Choice of tiers and positioning.}
The tiers differ in what they assume about normal images.
Tier~1 (template difference) assumes pixel alignment and is a lower bound.
Tier~2 (PaDiM) assumes one Gaussian per feature location; it is cheap, interpretable per location and
has few hyperparameters, so it is our primary method.
Tier~3 (PatchCore) drops the alignment assumption at the cost of a memory bank.
Unlike prior VisA work, which reports threshold-free scores on an official test set that is 55\%
anomalous, we fix thresholds on a labelled validation split and report operating points at a 1\% and
5\% deployment prevalence.
